% !TEX root = ../jlt.tex
\section{Partial Evaluation at Build Time}
\label{sec:binding}

Every function in this paper has run twice already.
The test behind the listings calls the same functions at runtime, on values it did not have at compile time.
Rule~1 is what makes that unremarkable: a \cppinline{constexpr} function of its arguments alone has no binding time of its own.\footnote{A \cppinline{consteval} function is evaluated only at compile time, so it has one binding time, which is why Rule~1 asks \cppinline{constexpr} and not \cppinline{consteval}.}
Theorem~\ref{thm:simulation} says the two evaluations agree, which is the presumption of this section: moving a call from compile time to run time changes where it is evaluated, not what it evaluates to.
Both are partial evaluation~\cite{jones1993partial} at build time: the compiler evaluates what it can and leaves the processor the rest.
Listing~\ref{lst:sites} is the paper's exhibit on both of its panels, each followed through the translation.
With the sequence free, the compiler evaluates what it can, the table, and keeps the fold over the sequence, which only the run time can finish.
What remains is the loop the optimiser emits, read as a function after Appel~\cite{appel1998ssa}: at \texttt{-O2} it allocates nothing and calls nothing, and each letter is four loads, the letter, its column, the next state and whether that state accepts; each call first copies the table by value, as Rule~1 has it.
With the sequence given and bound by \cppinline{constexpr}, the term evaluates to its count, and so does the code: Figure~\ref{fig:arrows}'s bottom edge on one program.
Given but not bound, the optimiser keeps the loop, so the constant is the binding's and not the optimiser's.

\begin{listing}[htbp]
\begin{minipage}[t]{0.49\linewidth}
\begin{cpp}
// Σ, the alphabet :
constexpr Alphabet bases = abc("ACGT");
// Σ* = ⊔ₙ Σⁿ :
constexpr Words oligos = words(bases);
// ~rev, the reverse complement: CCTA ↦ TAGG
// { w ∈ Σ⁶ | ~rev w = w }, the palindromic hexamers :
constexpr auto pali6 =
  oligos & len(6) & fixed_by(~rev);
// { (s, i) ∈ Σᴺ × ℕ | s[i, i+6) ∈ pali6 } :
constexpr auto cuts(size_t N) {
  return (oligos & len(N)) >> pali6;
}

static_assert(
  count(cuts(12)) == 1'835'008uz &&
  count(cuts(32)) == 27uz << 58);
\end{cpp}
\end{minipage}\hfill
\begin{minipage}[t]{0.49\linewidth}
\begin{cpp}
// those of pali6 beginning GA :
constexpr auto sites =
  pali6 & prefix("GA");

// the table, computed once by the compiler,
// a function of any s :
constexpr auto scanner =
  *automaton(sites);
static_assert(states(scanner) == 11uz);

// and only now a sequence :
constexpr auto dna = bytes(
  "GAATTCGAGCTCGGTACCCGGGGATCC"
  "TCTAGAGTCGACCTGCAGGCATGCAAGCTT");

// ten by definition, in windows of six :
static_assert(count(dna >> pali6) == 10uz);

// two by the table, in one pass :
// EcoRI's and SacI's
static_assert(
  count(dna >> scanner) == 2uz &&
  sites(bytes("GAATTC")) &&
  sites(bytes("GAGCTC")));
\end{cpp}
\end{minipage}
\begin{minipage}[t]{0.49\linewidth}
\begin{minted}{text}
-- the right panel's query at every N-mer, the
-- bind of the bag monad, counted by its laws:
-- a cylinder per offset, 4^3 · 4^(N−6), N − 5 alike
cuts N = Σ^N >>= (λ s. s >> pali6)
count (cuts N)
       = (N − 5) · 4^3 · 4^(N − 6)
count (cuts 32) ~> 27 · 2^58
         -- a constant; no sequence built
\end{minted}
\end{minipage}\hfill
\begin{minipage}[t]{0.49\linewidth}
\begin{minted}{text}
-- the table is a closed term,
-- the pass over the sequence is not
automaton sites ~> (0, T, F)
  -- static, 11 states: the compiler's
count (s >> scanner) =       -- the residual
  π₂ (fold (λ (q, c) a. (T q a, c + F (T q a)))
           (0, 0) (map rank s))
automaton pali6 ~> (0, T', F')
  -- 277 states: past the step limit, so
  -- the same function runs at run time
\end{minted}
\end{minipage}
\caption{Palindromic DNA restriction sites, counted two ways on one definition.
A palindrome is a word equal to its reverse complement, a $\theta$-palindrome for $\theta$ that antimorphic involution~\cite{kari2010watson}, as EcoRI's \texttt{GAATTC}; an enzyme cuts where one occurs.
Left: how many cuts a cocktail of every six-cutter makes in all $4^N$ $N$-mers, each digested once.
Right: the cuts of a smaller cocktail, the six-cutters beginning \texttt{GA}, in any one sequence, the caller's bytes read where they lie; here the polylinker of the plasmid pUC19.
Below each: its translation (Section~\ref{sec:translation}), a closed term the compiler reduces and a residual for the run time.
Without the count's laws, the left panel would build the $N$-mers and exhaust the step limit (Table~\ref{tab:binding}) before reaching its \cppinline{static_assert}.
The table is Brzozowski's automaton~\cite{brzozowski1964derivatives}, for one site the matcher of Knuth, Morris and Pratt~\cite{knuth1977fast}, which partial evaluation derives~\cite{consel1989partial}.
Its cost is moved to compile time where the step limit admits: ``negative overhead abstraction''~\cite{stroustrup2026peterman}.}
\label{lst:sites}
\end{listing}
The presumption holds where both evaluations finish and neither meets undefined behaviour.
Table~\ref{tab:binding} lists where the two binding times part.

\begin{table}[htbp]
\centering
\small
\begin{tabularx}{\linewidth}{@{}>{\raggedright\arraybackslash}p{0.18\linewidth}>{\raggedright\arraybackslash}X>{\raggedright\arraybackslash}X@{}}
\toprule
& at compile time, in constant evaluation & at run time \\
\midrule
who evaluates & the front end's interpreter, one step at a time & the processor, after the optimiser \\
a fold's length & at most 1\,048\,576 steps (\texttt{-fconstexpr-steps}), past which the program does not compile & the count alone \\
a call's depth & at most 512 nested calls (\texttt{-fconstexpr-depth}) & the stack \\
undefined behaviour & refused: the expression is not a constant, and the program does not compile & no diagnostic required and no guarantee: the optimiser may assume it never happens \\
a shortcut & only by resolution, a declared law, as Listing~\ref{lst:laws}'s overloads & the optimiser's, as it fuses Listing~\ref{lst:sites}'s stages into one loop \\
telling them apart & \texttt{if consteval} and \texttt{is\_constant\_evaluated()} say compile time & they say run time, and FJlt desugars it to the first where both agree \\
\bottomrule
\end{tabularx}
\caption{Where the two binding times differ, for Clang~23 with its default limits.
Each limit counts the evaluator's work, and can be raised.
Apart from the last row, no difference is a different value: where both evaluations finish without undefined behaviour, they agree (Theorem~\ref{thm:simulation}).}
\label{tab:binding}
\end{table}

\paragraph{The two evaluators are different programs.}
In constant evaluation no optimiser runs: the front end interprets a fold one step at a time, and the optimiser that fuses Listing~\ref{lst:sites}'s stages into one residual pass sees only the code emitted for run time.
So at compile time only a declared law shortens a fold, chosen by resolution, as for Listing~\ref{lst:laws}'s overloads: it buys at build time what the optimiser finds at run time by itself.

\paragraph{A law can leave the compiler a batch.}
Listing~\ref{lst:sites}'s section carries a length and a law in its type.
The length is a bound, as large as one likes but known before the scan, so the language is finite; with no bound, palindromes are not even regular~\cite{kari2010watson}, and no table holds them.
The law keeps the finite language small: the table lists what the law states, 277 states for one pair test per offset.
So one more overload of \cppinline{count} takes no table: a window is independent pairs, and each position its own lane.
Written as folds over small blocks, it vectorises at \texttt{-O2} with no intrinsic, about five times the table's speed (Appendix~\ref{app:fjlt:ssa}); the overload names a law of biochemistry and no instruction of any machine, and \Cpp{26}'s \cppinline{std::simd}~\cite{kretz2024simd} would make the lanes explicit.

\paragraph{The same laws license parallelism.}
A parallel \cppinline{concat} is one more overload, licensed by \cppinline{IsCommutativeMonoid}.
The branch on the binding time that selects it desugars to its compile-time branch, because the same law proves the branches agree (Appendix~\ref{app:fjlt:consteval}).

\Cpp{} passes no dictionaries but instantiates, so the translation models what a constrained call means, not what the compiler emits.
Instantiation is partial evaluation of the dictionary-passing term at a static type, so the two agree on every closed term, and parametricity holds of every instantiation (Appendix~\ref{app:fjlt:ssa}).

